% 第 37 章　理论光学的统一图景
\chapter{理论光学的统一图景}

\step{反常识开场}

第 33 章讲过针孔相机：一个不透光的盒子，一面扎个小孔，对面贴一张描图纸，窗外的树就在纸上倒立着显影。当时我们用“光沿直线传播”解释了一切：树上每一点向四面八方发光，只有穿过小孔的那一小锥能抵达屏上，物点对应像点，干净利落。

顺着这个解释往下推，你会得到一个斩钉截铁的预言：孔越小，每个物点被筛出的光锥越细，屏上的像点就越小，像就越锐利。推到极致——把孔缩成一个数学上的点——每个物点只放出一条光线，像应该完美清晰，只是暗一些。“孔越小越清晰”，这几乎是从“光走直线”直接流出来的结论，不容置疑。

可任何真动手做过针孔相机的人都知道结果：孔扎到针尖那么细的时候，画面不但没有变得刀锋般锐利，反而整个软掉、糊掉，像隔着一层毛玻璃。“光走直线”这位忠实的老仆人，在它自己的极限处背叛了自己。

同样憋着一口气的还有三个日常现象。手机摄像头动辄五千万像素，可把小字推近再推近，细节总有一个擦不掉的天花板；光学显微镜再贵，也看不清病毒的形状——病毒的直径大约一百纳米，比可见光的波长还小；调频收音机在大楼背后照样响，可见光却被一面墙挡得严严实实——都是波，凭什么待遇天差地别？

这四个问题，靠“射线”或“波”任何一方单独出面都答不全。本章的任务，是把前几章看似两本教材的射线光学与波动光学熔成一张地图：它们不是两种光，而是同一套方程在不同放大倍率下的两张照片。

\step{先猜一猜}

在继续之前，请写下你对“小孔反而更糊”的猜测。三种假说，哪一种最像真相？

第一种，工艺说：孔边缘的毛刺和厚度把光碰散了，换成激光切割的精密圆孔就不会糊。

第二种，妥协说：有两股方向相反的势力在较劲——孔大时这一股占上风，孔小时那一股占上风，中间某处存在一个“最佳孔径”，糊与锐在那里交接。

第三种，波动说：光根本不完全走直线，“直线”只是一种近似，而且孔越小，这个近似越糟糕。

光写假说还不够，请再押一个数：如果真有最佳孔径，它大约是头发丝的粗细（约 \(0.05\) 毫米）、缝衣针的眼（约 \(0.5\) 毫米），还是铅笔芯（约 \(2\) 毫米）？把你的两票都写在纸上。本章结尾，我们不仅要对答案，还要把最佳孔径当场算出来——如果你的直觉押对了数量级，那说明你已经摸到了本章最硬的公式。

\step{动手实验}

\begin{experiment}[三个孔与墙上的一圈光环]
材料：一张厚黑卡纸或铝箔、一枚缝衣针、牙签、胶带、两个能套在一起的纸筒（或一个鞋盒）、一小片描图纸或蜡纸、一支激光笔。

第一步：做一台三孔针孔相机。在卡纸上并排扎三个孔：用牙签捅一个约 \(2\) 毫米的粗孔；用针尖正常扎一个约 \(0.5\) 毫米的中孔；再用针尖轻轻抵住纸面旋转，扎一个刚能透光的微孔（约 \(0.1\) 到 \(0.2\) 毫米）。把卡纸蒙在纸筒一端，另一端蒙上描图纸。

第二步：白天把纸筒对准窗外明亮的景物（台灯也可以），轮流只留一个孔透光，比较三幅像。你会看到：粗孔的像亮但糊成一团；中孔的像最锐利；微孔的像暗而软，边缘发毛。把三幅像的“亮度”和“锐度”分别打分记在纸上——“最清晰的像来自中间的孔”，这就是被“孔越小越清晰”的预言漏掉的事实。

第三步：进暗室，把激光笔对准那个最小的孔，让透过的光投到两三米外的白墙上。墙上不再是一个红点，而是一枚亮斑套着一圈一圈同心圆环的图案，孔越小，整个图案越大。换粗孔重做，环缩小成几乎一个点。请盯着光环看足一分钟：沿直线飞行的东西，凭什么在孔外画出一圈一圈的环？
\end{experiment}

三个观察请带走：存在最佳孔径；微孔给出的是“图案”而不是“像”；图案的尺寸与孔径反着走。接下来看旧模型怎么接招。

\step{旧模型遇到麻烦}

先让射线模型（第 33 章）答辩。它的预言是单调的：孔径 \(d\) 缩小，每个物点在屏上的几何模糊斑就正比于 \(d\) 缩小，锐度只升不降。这个模型里根本没有“最佳孔径”这个词条，更没有“光环”——沿直线飞行的东西，凭什么绕过孔的边缘向外拐弯？射线模型可以解释粗孔的糊，却对微孔的糊束手无策。它没全错，但它的管辖范围显然有个边界，而我们刚刚摸到了那条边界。

再让波动模型（第 35 章）答辩。它手里有衍射，知道波过孔会散开，光环正是散开后的干涉图案。但它也有三处窘迫。第一，它说不清“直线传播”这位老朋友是从哪条定理里掉出来的：声波也是波，凭什么声波绕柱而行、光波偏要直走？第二，第 34 章的透镜成像在波的语言里还没有着落：一块玻璃凭什么把散开的波重新收拢成一个点？第三，两支手电筒的光叠在一起从来不出现条纹，为什么有的波配干涉、有的波不配？波动模型有一口袋现象，却缺一张总图。

两个模型背后有没有同一个老大？有，而且第 25 章已经请他出过场：麦克斯韦方程。光是电磁场的一种运动方式，而电磁场的一举一动——它的空间分布和它的随时间变化——共同服从麦克斯韦方程组，波只是这组方程天然携带的解。本章剩下的内容，就是从这同一个出发点，把射线、衍射、成像、相干性逐一指认为它的推论。

\begin{mainline}
统一图景可以压成一条算式：\textbf{一个对象（电磁场）＋一组方程（麦克斯韦方程）＋边界条件（孔、缝、透镜、屏）＝全部光学现象。}射线不是被推翻，而是被定位成“波长趋于零”的极限情形；衍射是波前被边界裁剪后的必然变形；成像是人类把边界条件精心设计成透镜的结果。
\end{mainline}

\step{发明新概念}

要从方程走到现象，我们需要四件新工具：波前、相位、相干性和空间频率。前三件本节造出来，第四件留给数学翻译器。

\textbf{波前与相位}。波场里的每一点都随身带着两个数：振幅（此刻最多能振多大）和相位（此刻振到节拍的哪一格）。相位像全班大合唱时每人手里的节拍器读数：读数相同的点连成一片，就是\textbf{波前}。点光源的波前是一层层同心球面；离光源足够远，球面的一小片看起来就是平的，得到平面波前。它像梯田的等高线——等高线是“高度相等”的线，波前是“相位相等”的面，顺着垂直方向走变化最快；它不像等高线的地方更要记住——山不会跑，波前却以光速向外推进，而且相位不是高度，它是个循环的量，转满一圈就归零。

现在给出本章最重要的定义：\textbf{光线就是波前的垂线}。射线不是光的本体，而是从波前派生出来的记号——波前朝哪边推进，就在那个方向上画一条线。平面波前的垂线彼此平行，所以均匀空间里的“光线”是直线；点光源球面波前的垂线向外辐射，所以影子边缘锐利。“光走直线”从此有了出处：它不是公理，是波前几何的推论。

\begin{center}
\begin{tikzpicture}[scale=1.0,>=Stealth]
  \fill (0,0) circle (1.6pt) node[below=2pt] {点光源};
  \foreach \r in {0.7,1.4,2.1} {\draw (0,0) circle (\r);}
  \foreach \a in {15,75,135,195,255,315} {\draw[->] (\a:2.1) -- (\a:2.95);}
  \node at (135:1.8) [above left] {波前（相位相等的面）};
  \node at (15:3.1) [right] {光线（波前的垂线）};
\end{tikzpicture}
\end{center}

第 35 章用过的惠更斯作图法，现在升级为正式工具：波前上每一点都当作新的小波源，发出半球形的子波，下一刻的波前就是所有子波共同的外包络。直线传播、反射、折射、绕孔散开，全都能用这把圆规直尺作出来。它不是额外的新假设——在“相位在一个波长内变化远大于振幅变化”的条件下，可以从麦克斯韦方程把它推出来，推导的草图放在挑战层。

\textbf{相干性与激光}。两支手电筒照同一面墙，为什么永远没有条纹？因为普通光源里亿万个原子各自独立发光，每过极短一瞬间，整束光的相位就重新洗牌一次，干涉条纹的位置每秒狂跳千万亿次，任何眼睛和底片都只能收到洗匀后的平均值。用本章的语言说：普通光源的\textbf{相干长度}极短，只有微米量级，两束光必须几乎同步到达才能“认亲”。

激光则是另一个物种，本章先只介绍现象，不追根源。拿激光笔照白墙，光斑不是均匀的红点，而是细密的颗粒状明暗，头稍微一动，颗粒就滚动游走——这叫\textbf{散斑}：墙面凹凸不平，反射出的无数子波带着稳定的相位关系在你眼前互相干涉。普通台灯照不出散斑，因为相位队伍太乱。激光的相位队伍整齐得反常，相干长度可达几十厘米甚至更远，所以干涉、衍射的一切效应都被它放大到肉眼可见。为什么激光能把相位排成队？那要翻开原子发光的微观账本，是第 55 章“光与原子的合作”的课题。本章只需要记住现象：存在相位高度整齐的波，而相干性决定了“什么样的波才配稳定干涉”。

\step{一句话模型}

\begin{oneline}
光是满足麦克斯韦方程的电磁场；边界条件（孔、缝、透镜）决定波前如何展开；所谓“光线”只是波长足够小时波前的垂线，衍射是波前被裁剪后的必然变形，成像则是透镜对各个空间频率成分的先分解、再合成。
\end{oneline}

这句话请拆开读两遍。第一遍读“波长足够小”：射线光学没有死，它被收编为一个极限，本章的每条公式都应该在波长趋于零时退回第 33 章。第二遍读“边界决定展开”：针孔、狭缝、透镜、光栅，本质上都是给波前做手术的不同刀法；光学工程师的全部手艺，就是设计刀法。

\step{数学翻译器}

直觉到位，开始翻译。第一道工序：把“波”写成式子。第 25 章从麦克斯韦方程得到过波动方程，一维情形下它长这样：
\[
  \frac{\partial^2 E}{\partial x^2}=\frac{1}{c^2}\frac{\partial^2 E}{\partial t^2}.
\]
左边是场在空间上的弯曲程度，右边是它随时间变化的快慢，光速 \(c\) 是两者之间唯一的兑换率。它有一个干净的解：
\[
  E=E_0\cos(kx-\omega t),
\]
其中 \(k=2\pi/\lambda\) 叫波数（单位长度里塞几个波长），\(\omega=2\pi f\) 叫角频率（单位时间里转几格相位），并且 \(\omega/k=c\)。

按规矩，立刻用特殊情形检查。盯住空间某一点（令 \(x=0\)）：\(E=E_0\cos(-\omega t)\)，一个原地上下振荡的量，频率正是 \(f\)——合理。按下暂停键（令 \(t=0\)）：\(E=E_0\cos kx\)，空间上一条余弦曲线，周期正是 \(\lambda\)——合理。盯住某个波峰（令 \(kx-\omega t=\text{常数}\)）：解出 \(x=(\omega/k)t=ct\)，波峰以光速向右跑——合理。把括号里换成 \(kx+\omega t\)，波峰就向左跑：方向反转检查通过。一个式子，把时间振荡、空间周期、传播方向和速度全部说清，这就是波动方程的威力。

第二道工序：衍射角。波前被宽度为 \(D\) 的孔裁掉两侧之后，不再平，而是向一个角度范围散开。用惠更斯作图可以证明，单缝后方的第一暗纹出现在
\[
  \sin\theta=\frac{\lambda}{D}
\]
的方向上（圆孔对应的角度是 \(1.22\lambda/D\)，系数略有差别，数量级相同，本章统一按 \(\theta\approx\lambda/D\) 估算）。它回答“孔把光撒得多开”：\(\lambda\) 是波的本性，\(D\) 是你给波前的裁剪尺度。

\begin{center}
\begin{tikzpicture}[scale=1.0,>=Stealth]
  \draw[very thick] (0,1.0) -- (0,2.6);
  \draw[very thick] (0,-1.0) -- (0,-2.6);
  \draw[<->] (-0.18,-1.0) -- (-0.18,1.0) node[midway,left] {$D$};
  \draw[->] (0,1.0) -- (4.0,2.6);
  \draw[->] (0,-1.0) -- (4.0,2.6);
  \draw[dashed] (0,1.0) -- (2.2,1.0);
  \draw (1.6,1.0) arc[start angle=0,end angle=27,radius=1.6];
  \node at (14:2.1) {$\theta$};
  \draw[very thick] (4.0,0.4) -- (4.0,3.4);
  \node at (4.0,3.6) {远处的屏};
  \node at (2.0,-1.8) {边缘两束子波的光程差达到 $\lambda$ 时，方向 $\theta$ 上相消};
\end{tikzpicture}
\end{center}

检查极限：\(\lambda\) 趋于零时 \(\theta\) 趋于零，波乖乖直走——射线极限回归；\(D\) 趋于无穷大时 \(\theta\) 趋于零，孔大到等于没剪，不散；\(D\) 缩到与波长差不多时 \(\theta\) 接近一个弧度量级，“孔”已经谈不上是孔，而是一个向四面八方辐射的点源。三个方向都合理。

第三道工序：针孔的最佳孔径，本章的招牌计算。孔径为 \(d\)、针孔到屏的距离为 \(L\)、景物在远处时，两股势力分别是：几何模糊斑，大小约为 \(d\)（孔越大，每个物点在屏上糊开的斑越大）；衍射模糊斑，大小约为 \(\lambda L/d\)（孔越小，波散得越开，\(L\) 越长散得越宽）。总模糊大致是两股相加：
\[
  \text{模糊}\approx d+\frac{\lambda L}{d}.
\]
第一项随 \(d\) 增大，第二项随 \(d\) 减小，两者相等处总模糊最小，于是最佳孔径
\[
  d^{*}\approx\sqrt{\lambda L}.
\]
代入真实数据：取可见光波长 \(\lambda\approx 550\) 纳米、纸筒长 \(L=20\) 厘米，算出 \(d^{*}\approx\sqrt{1.1\times10^{-7}\ \mathrm{m}^2}\approx 0.33\) 毫米。这正是动手实验里“中孔最清晰”的位置，也正是你猜测单上“缝衣针的眼”那一档。直觉、实验与公式三方会师。

\begin{center}
\begin{tikzpicture}[scale=1.35,>=Stealth]
  \draw[->] (0,0) -- (3.6,0) node[right] {孔径 $d$};
  \draw[->] (0,0) -- (0,2.4) node[above] {模糊斑};
  \draw[domain=0.05:3.2,smooth] plot (\x,{0.55*\x}) node[right] {几何 $\sim d$};
  \draw[domain=0.35:3.2,smooth] plot (\x,{0.62/\x}) node[right] {衍射 $\sim\lambda L/d$};
  \draw[dashed,domain=0.32:3.2,smooth] plot (\x,{0.55*\x+0.62/\x});
  \fill (1.06,1.17) circle (1.4pt);
  \draw[dotted] (1.06,0) node[below] {$d^{*}$} -- (1.06,1.17);
\end{tikzpicture}
\end{center}

检查极限：\(\lambda\) 趋于零时 \(d^{*}\) 趋于零，“孔越小越好”——射线模型的单调预言完好无损地回来了，两个模型在极限处握手。\(L\) 越大 \(d^{*}\) 越大：长筒相机的孔要开得稍大些，也合理。

顺手把人眼也算一遍。瞳孔直径白天约 \(2\) 毫米、夜里约 \(7\) 毫米，取中间值 \(4\) 毫米：衍射角 \(\theta\approx\lambda/D\approx 550\ \mathrm{nm}/4\ \mathrm{mm}\approx1.4\times10^{-4}\) 弧度，约合半角分。而视力表 \(1.0\) 那一行的缺口，张角正好约一角分。视网膜上感光细胞的间距刚好够接住衍射送来的信息，一分不多一分不少——眼睛的分辨率尽头不是细胞不够好，而是波的本性画下的线。

\begin{mathbridge}
相位的记账最好交给复数。把场写成复振幅 \(A\mathrm e^{\mathrm i\varphi}\)（欧拉公式 \(\mathrm e^{\mathrm i\varphi}=\cos\varphi+\mathrm i\sin\varphi\)，真实的场取实部），相位变成复平面上的转角，叠加变成复数相加，强度变成模的平方。双缝干涉在这个语言里只剩一行：两束复振幅相加再取模方，交叉项自动长出 \(\cos\Delta\varphi\)，条纹明暗全部到位。虚数单位 \(\mathrm i\) 在这里没有任何神秘含义，它只是“比另一束晚四分之一个节拍”的缩写。
\end{mathbridge}

\begin{mathbridge}
空间频率是本章第四件工具。把一幅图像沿空间方向“听”成和弦：大片平缓的明暗是低音（低空间频率），细密的纹理和锐利的边缘是高音（高空间频率，单位用“每毫米多少线对”）。远场衍射理论给出一个惊人的结论：\textbf{孔径后面的远场图案，就是孔径形状的傅里叶变换}——方向即空间频率。而透镜把方向变成位置（第 35 章的直觉），所以透镜的焦平面就是频谱面：一束平行光（一个空间方向）聚成焦面上一个点。于是成像这个动作可以重新描述为：透镜先把物分解成各个空间频率，再在像面上重新合成。但透镜口径有限，只能接住角度不太大的成分，太高的空间频率直接流失——\textbf{任何成像系统都是一台低通滤波器}，像永远比物“钝”一点。阿贝在一八七三年用这套语言解释了显微镜：携带病毒一百纳米细节的空间频率，对应的角度大到镜头根本接不住，于是那些细节从来没能参加合成。这套理论至今仍是光刻机与显微镜设计的宪法。
\end{mathbridge}

\begin{challenge}
近场与远场由菲涅耳数 \(F=D^2/(\lambda L)\) 划分。\(F\gg1\) 时是近场：屏上的图案是孔的几何影子镶一圈毛边；\(F\ll1\) 时是远场：图案长成孔径的傅里叶变换，此后距离再远也只是等比放大，不再变形。激光笔光斑约 \(2\) 毫米、波长 \(650\) 纳米，过渡距离 \(D^2/\lambda\) 约 \(6\) 米——你桌上的激光实验恰好骑在近场与远场的分界线上。更深一层：携带比波长更细结构的那些空间频率，沿传播方向不是振荡而是指数衰减（称为倏逝波），它们根本不出门。所以衍射极限的真正根源不是“波会散开”，而是“细节不旅行”。近场光学显微镜把一根细探针贴到样品几十纳米之内，在倏逝波死光之前把它截获——衍射极限就这样被绕过（注意，是绕过，不是打破）。
\end{challenge}

\begin{challenge}
惠更斯原理可以严格化。基尔霍夫证明：屏后任意一点的场，等于孔面上各点的场及其变化率的加权和，权重里带一个倾斜因子——这正是子波只向前包络、不向后倒流的原因。而射线光学也能从波动方程里逼出来：当相位在一个波长内的变化远大于振幅的变化时，波动方程退化成一条关于相位的方程（程函方程），它的推论正是直线传播、反射定律与折射定律。几何光学三定律，一条不剩，全是“波长趋于零”的推论。
\end{challenge}

\step{模型的边界}

统一图景再威风，也有它的边界，逐条列清。

射线近似的边界：只在“器件尺度远大于波长”时好用。大楼对调频广播（波长米级）只是一条窄缝，波从容绕过；同一栋楼对可见光（波长约半微米）宽如十个足球场，射线近似好得惊人，影子刀切般锐利。所以收音机穿楼而光不穿楼，不是两种波“性格不同”，而是同一条 \(\theta\approx\lambda/D\) 在两种尺度下的两副面孔——这是一个专门用来绊倒直觉的反例，值得记牢。

标量近似的边界：本章把场当成一个数来算，在近轴、小角度时很好；强聚焦、大角度时必须请回偏振（第 36 章），把场当矢量处理。

衍射极限的边界：它只对“自由传播到远场”成立。近场探测、特殊设计的透镜能绕过它，但绕过的依据仍是同一套麦克斯韦方程——方程没有被违反，被绕开的只是“自由传播”这个前提。

相干性的边界：条纹要稳定，相位差必须稳。光源面积太大、颜色太杂，都会把条纹洗没；散斑则是相干性的勋章兼麻烦——激光投影仪里要专门加一片消散斑的毛玻璃。

经典场图像的天花板：把光调暗再调暗，暗到极致，探测器上出现的不是均匀变淡的辉光，而是一颗一颗分立的“咔哒”。连续摊开的经典场预言错了。这不是麦克斯韦方程的笔误，而是它的边界：光的能量交换另有一本量子账，从第 42 章“经典物理遇到的几场灾难”开始清算。同样，激光“为什么”相干，本章只押不付，微观机制留给第 55 章。

\step{回头解释开场现象}

开场四案，逐一销账。

针孔相机三孔：粗孔时几何模糊（正比于 \(d\)）当家，像糊；微孔时衍射模糊（正比于 \(\lambda L/d\)）当家，像软；中孔落在 \(d^{*}\approx\sqrt{\lambda L}\approx 0.3\) 毫米附近，两股势力恰好交接，像最锐。你的猜测对答案：妥协说与波动说本是同一件事的两半，工艺说出局——孔缘毛刺只添一点噪声，主力军是波自己。墙上那圈同心环，则是小孔的远场衍射图样：圆孔的傅里叶变换是一枚中心亮斑套一圈圈暗亮相间的环，孔越小，图案越大，与实验完全对上。

手机像素的天花板：镜头口径 \(D\) 小，衍射斑（约 \(\lambda f/D\)）一落下来就盖住好几个像素，像素再多，也只是把同一枚模糊斑数得更细。厂商拼命“加大底”，加的不是像素，是 \(D\)。

显微镜看不见病毒：携带百纳米细节的空间频率对应的角度超出镜头接收范围，细节从未参加合成。电子显微镜的对策是换用波长只有可见光几万分之一倍的物质波——为什么电子也是波？那是第 42、43 章的开场戏。

收音机穿楼：波长米级对楼宽，\(\lambda/D\) 巨大，“影子”根本来不及形成。光与无线电没有两副心肠，只有一把尺子。

\step{脑洞实验与挑战题}

\begin{thought}[波长一米的世界与地球那么大的望远镜]
做两个方向相反的尺度游戏。第一个：假如可见光的波长变成一米，其余一切不变。你的瞳孔远小于波长，每只眼睛退化成一枚只能测亮度的点探测器，成像、影子、针孔相机全部作废，“光走直线”从日常词典里消失，人类文明大概会先发明波动光学。我们恰好住在波长远远小于日常尺度的世界，所以童年物理从射线开始——射线模型是环境送给我们的礼物，不是光的真面目。

第二个方向反过来：把口径推到地球那么大。这不是幻想，二〇一九年人类真的干过：分布全球的八台射电望远镜联网同步观测，等效口径约一万公里，在波长 \(1.3\) 毫米处工作，角分辨率 \(\theta\approx\lambda/D\approx1.3\times10^{-10}\) 弧度，约二十五微角秒——正好数清 M87 星系中心黑洞“影子”的轮廓（约四十微角秒）。衍射极限在这里不是牢笼，而是一封邀请函：它邀请你把全世界的天线粘成一台仪器。至于那个被拍下来的影子为什么会弯曲成环，请翻到第 41 章的弯曲时空。
\end{thought}

\begin{puzzles}
\item 看不清黑板时眯起眼会清楚些：眼睑把瞳孔挤小，减少了离焦造成的几何模糊。但眯到极致反而更糊。请解释为什么存在“最佳眯缝宽度”，并估计它的量级。（提示：把眼皮当成可调的光阑，直接套用针孔相机的权衡公式 \(d^{*}\approx\sqrt{\lambda L}\)，想想这里的 \(L\) 是眼球的长度。）
\item 同一台针孔相机，分别用红光和蓝光拍摄，哪一张照片更锐利？换颜色时最佳孔径要不要跟着换？（提示：几何模糊与颜色无关，衍射模糊正比于波长；\(d^{*}\approx\sqrt{\lambda L}\)，红光的最佳孔径更大。）
\item 射电望远镜 FAST 的口径有五百米，一台普通一米口径的光学望远镜和它比，谁的角分辨率更高？大约差多少个数量级？（提示：比的不是口径而是 \(\lambda/D\)。射电波长取 \(0.1\) 米量级，可见光取 \(5\times10^{-7}\) 米；口径大了五百倍，波长却大了几十万倍。）
\item 激光照白墙看到的散斑，颗粒的大小由什么决定？戴上眼镜和摘掉眼镜看，斑点会跟着变形吗？（提示：散斑不在墙上，而在你眼前的整个空间里；颗粒大小由接收它的瞳孔的衍射决定，约等于瞳孔给出的衍射斑。）
\end{puzzles}
